\section{preliminary}

% 介绍推理引擎的执行模式
\subsection{Inference System Execution Modes}
\label{sec:3.1}
% 介绍推理引擎的普通范式
Modern LLM inference systems execute requests iteratively with continuous batching. At each iteration, the \textbf{scheduler} updates the set of running requests and constructs the execution configuration for the current batch. 
The resulting scheduling information is then dispatched to \textbf{workers}, which perform model execution \textbf{\texttt{Execute\_Model}} followed by token processing \textbf{\texttt{Sample\_Tokens}}. 
We therefore abstract each inference iteration as:
\[
\text{Scheduler}
\rightarrow
\text{Workers}\left(
\texttt{Execute\_Model}
\rightarrow
\texttt{Sample\_Tokens}
\right).
\]
% 介绍 Non-SD 和 SD 两种模式的区别
\begin{itemize}
    % Non-SD
    \item \textbf{Non-SD execution mode}. \texttt{Execute\_Model} performs a standard LLM forward pass using a \textbf{decode kernel} and produces logits for the next token of each request.
    \texttt{Sample\_Tokens} then samples the next token from the resulting logits.
    % SD
    \item \textbf{SD execution mode}. \texttt{Execute\_Model} performs LLM verification over multiple draft tokens in parallel using a \textbf{prefill-style kernel} for multi-token queries. \texttt{Sample\_Tokens} then invokes the drafter to generate draft tokens for subsequent verification.
\end{itemize}
% 说明将二者强行统一成一条路径的问题
Consequently, Non-SD and SD differ in both execution flow and kernel usage. Unifying them into a single execution path would require extra alignment or force one mode to use a suboptimal kernel, introducing unnecessary overhead.

% Agent 任务的调度优化问题定义
\subsection{Scheduling Objective for Agent Tasks}

% 简要介绍 agent task 的特点
Unlike conventional LLM tasks, agent tasks are typically executed as multi-step workflows consisting of multiple dependent steps. 
For all serial agent workflows considered in this work, we abstract their execution as a ReAct-style pattern, where LLM inference and tool execution alternate throughout the workflow~\citep{ReAct}. 
As a result, the end-to-end performance of an agent task depends on the execution of the entire workflow rather than the latency of any individual LLM request.

% 每个 agent 请求 R_i 包括 K_i 个 Step，每个 Step 包括一次 LLM 执行和一次 Tool 执行。
For an agent request $R_i$, we represent its workflow as an ordered
sequence of dependent steps,
\begin{equation}
R_i = \langle S_i^1, S_i^2, \ldots, S_i^{K_i} \rangle,
\end{equation}
where $K_i$ denotes the total number of steps of $R_i$.
Each step $S_i^{k}$ consists of an LLM execution followed by an optional tool
execution before the next step can be issued.

% 一个 agent 请求 R_i 在某种调度策略下的总端到端 JCT 的构成
Let $T_{\pi}^{\mathrm{Queue}}(S_i^j)$ denote the waiting time of step $S_i^j$ under scheduling policy $\pi$, $T^{\mathrm{LLM}}(S_i^j)$ denote its LLM execution time, and $T^{\mathrm{Tool}}(S_i^j)$ denote the tool execution time. The end-to-end Job Completion Time (JCT) of the agent request $R_i$ can be expressed as
\begin{equation}
\mathrm{JCT}_{\pi}(R_i) = \sum_{j=1}^{K_i} \left( T_{\pi}^{\mathrm{Queue}}(S_i^j) + T^{\mathrm{LLM}}(S_i^j) \right) + \sum_{j=1}^{K_i-1} T^{\mathrm{Tool}}(S_i^j).
\end{equation}

% 优化问题：最小化一组 Agent 任务的平均端到端 JCT
Given a set of agent requests $\mathcal{R}$, our scheduling objective is to find a feasible scheduling policy $\pi$ that minimizes the average end-to-end JCT:
\begin{equation}
\pi^{*} = \arg\min_{\pi \in \Pi} \frac{1}{|\mathcal{R}|} \sum_{R_i \in \mathcal{R}} \mathrm{JCT}_{\pi}(R_i),
\end{equation}
where $\Pi$ denotes the set of feasible scheduling policies that respect the dependency between agent steps and the serving capacity of the inference system.

% 引用论文说明这个优化问题是 NP-Hard
This global scheduling objective is computationally challenging due to the dependencies among agent steps and the contention among concurrent requests. 
Astraea~\citep{Astraea} formalizes a similar multi-stage agent scheduling problem and has proved it to be NP-hard. 
Therefore, practical online agent serving systems typically rely on efficient   euristic policies to approximate the optimal schedule.